%% kapa: exam=1 round=25 year=2014 period=2 form=A subject=appraisal_law subject_ko=부동산관계법규 seq=27 number=27
%% source: Q-Net 감정평가사 기출문제 (cst00309) artlSeq=5101207 file=extracted/1008/5101207/2014년 제25회 감정평가사 1차시험 기출문제/2교시(A형).hwp
%% conversion: kordoc 4.24.51 -> tools/gen1.py (auto); review=proofread-sonnet; flags=none
\begin{question}{27}
국유재산법령상 국유재산의 구분과 종류에 관한 설명으로 옳은 것은?
\begin{choices}
  \item 국유재산은 그 형상에 따라 행정재산과 공공재산 및 일반재산으로 구분한다.
  \item 대통령 관저는 공용재산이 아니다.
  \item 국가가 비상근무에 종사하는 공무원에게 제공하는, 해당 근무지의 구내 또는 이와 인접한 장소에 설치된 주거용 시설은 공공용재산이다.
  \item 정부기업이 인사명령에 의하여 지역을 순환하여 근무하는 소속 직원의 주거용으로 사용하는 재산은 기업용재산이다.
  \item 국가가 보존할 필요가 있다고 국토교통부장관이 결정한 재산은 보존용재산이다.
\end{choices}
\end{question}
\answer{4}
